\section{Every finite original projective family retains its actual Segre line and sections}
\label{sec:finite-segre-original-projective-maps}
\begin{definition}[The actual finite Segre index and original morphism]
\label{def:finite-original-segre-morphism}
\lean{EnvelopeProperty.FiniteSegreOriginalProjectiveMaps.dimension,EnvelopeProperty.FiniteSegreOriginalProjectiveMaps.index,EnvelopeProperty.FiniteSegreOriginalProjectiveMaps.index_bijective,EnvelopeProperty.FiniteSegreOriginalProjectiveMaps.presentation,EnvelopeProperty.FiniteSegreOriginalProjectiveMaps.morphism,EnvelopeProperty.FiniteSegreOriginalProjectiveMaps.morphism_structure,EnvelopeProperty.FiniteSegreOriginalProjectiveMaps.morphism_isClosedImmersion}
\leanok
\uses{def:original-pair-segre-morphism,thm:original-segre-composite-closed-embedding}
Let $q_j:X\to\mathbf P_k^{d_j}$ be a nonempty finite family of actual
original morphisms with the same original base triangle. Iterating the
genuine binary Segre morphism constructs one actual original projective
morphism on the unchanged $X$. Its coordinates are in bijection with
all tuples of original factor coordinates. Its entire original base
triangle is retained. If $q_0$ is a closed immersion, this morphism is
a closed immersion.
\end{definition}
\begin{proof}
For one factor use the original map and its actual coordinate index.
For a larger family separate the first factor and the remaining nonempty
family. Apply the recursively constructed map to the remaining family,
and form its actual binary Segre composite with the first factor.
Both base composites are the original base, so the whole pair exists
and retains that base triangle. The binary coordinate-index bijection,
together with the recursively constructed tail bijection, gives a
bijection from all factor-coordinate tuples to all target coordinates:
its inverse first separates the binary index and then recursively
recovers every tail coordinate. If the first original map is a closed
immersion, the actual binary Segre closed-embedding theorem applies
irrespective of the remaining map. This proves the closed-embedding
assertion at every family length.
\end{proof}
\begin{theorem}[The entire original finite tensor line and every coordinate section]
\label{thm:finite-original-segre-whole-line}
\lean{EnvelopeProperty.FiniteSegreOriginalProjectiveMaps.tensorLine,EnvelopeProperty.FiniteSegreOriginalProjectiveMaps.tensorLine_isInvertible,EnvelopeProperty.FiniteSegreOriginalProjectiveMaps.tensorCoordinate,EnvelopeProperty.FiniteSegreOriginalProjectiveMaps.lineIso,EnvelopeProperty.FiniteSegreOriginalProjectiveMaps.lineIso_coordinate}
\leanok
\uses{def:finite-original-segre-morphism,thm:original-segre-pulled-line-iso,thm:original-segre-pulled-coordinate-section,lem:fc-original-tensor-section-homs}
The actual pullback of the target $\mathcal O(1)$ is isomorphic to the
entire original tensor product $\bigotimes_jq_j^*\mathcal O(1)$.
Under this derived whole isomorphism, every actual target coordinate
section is the tensor of the same tuple of original coordinate sections.
All these actual original line modules are invertible.
\end{theorem}
\begin{proof}
For one factor the comparison is the identity. For a larger family use
the genuine binary Segre line isomorphism, then tensor the first-factor
identity with the recursively constructed whole comparison for the tail.
Original projective pullbacks are invertible, and the tensor of actual
invertible modules is invertible. The binary comparison sends each
coordinate section to the two original factor sections. Naturality of
the actual tensor-section morphism then composes its second section
with the recursive comparison, whose coordinate identity gives the
entire original tail tensor. This proves the assertion for every tuple.
\end{proof}
\begin{theorem}[The full original section spaces and weighted families are isometric]
\label{thm:finite-original-segre-whole-section-isometry}
\lean{EnvelopeProperty.FiniteSegreOriginalProjectiveMaps.wholeSectionsEquiv,EnvelopeProperty.FiniteSegreOriginalProjectiveMaps.wholeSectionsNorm,EnvelopeProperty.FiniteSegreOriginalProjectiveMaps.weightedNormEval,EnvelopeProperty.FiniteSegreOriginalProjectiveMaps.coordinateNorm}
\leanok
\uses{thm:finite-original-segre-whole-line,thm:scheme-global-section-model-norm-isometry,thm:scheme-global-section-weighted-norm-isometry,thm:original-segre-coordinate-model-norm}
The derived whole line isomorphism gives an equivalence on the entire
original global section spaces. If $X$ is proper, it preserves every
actual proper-model section norm and every finite real-weighted norm
evaluation, with the same denominator and weights, at every native valued
point over a bounded base valuation. Every actual target coordinate norm
is the product of the corresponding original factor-coordinate norms,
including zero values and nongeneric supports.
\end{theorem}
\begin{proof}
Apply whole global sections to the constructed line isomorphism and its
actual inverse. Whole-module isometries preserve every section norm at
every actual proper valued point. Mapping a finite family through this
same equivalence leaves each logarithmic term and its real weight
unchanged, hence leaves its finite supremum and denominator-normalized
value unchanged, including empty and common-zero families. The binary
coordinate norm theorem identifies the first factor times the tail
coordinate norm. Repeatedly apply this identity to the tail. The finite
product over the entire family is exactly the first factor times that
tail product, proving the stated formula without nonvanishing conditions.
\end{proof}
These statements concern the actual finite original Segre embedding and
its whole original tensor line. Identifying its graded algebra with the
positive diagonal of the reconstructed multigraded model is a further
step; whole projective numerical-class finite dimensionality and the
paper topology are separate unresolved obligations.
